\clearpage
\section{선별 상세 풀이 I}
\label{sec:norm-trace-solutions-a}

\subsection*{연습문제 \ref{ex:norm-trace-quadratic-matrix}}

\begin{solution}
기저를 $\mathcal B=(1,\sqrt5)$라 하자. 먼저
\[
  \alpha\cdot1=3+2\sqrt5,
\]
\[
  \alpha\cdot\sqrt5
  =(3+2\sqrt5)\sqrt5=10+3\sqrt5.
\]
따라서 열벡터 관례에서
\[
  [m_\alpha]_{\mathcal B}
  =\begin{pmatrix}
    3&10\\
    2&3
  \end{pmatrix}.
\]
그러므로
\[
  \Tr_{L/\Q}(\alpha)=3+3=6
\]
이고
\[
  \operatorname N_{L/\Q}(\alpha)
  =3\cdot3-10\cdot2=-11.
\]

이차확장의 켤레를 사용하면 같은 결과를 더 빠르게 얻는다.
\[
  \overline\alpha=3-2\sqrt5,
\]
\[
  \Tr(\alpha)=\alpha+\overline\alpha=6,
  \qquad
  \operatorname N(\alpha)=\alpha\overline\alpha=9-20=-11.
\]

역원은
\[
  \alpha^{-1}
  =\frac{\overline\alpha}{\alpha\overline\alpha}
  =\frac{3-2\sqrt5}{-11}
  =\frac{2\sqrt5-3}{11}.
\]
노름의 곱셈성에서
\[
  \operatorname N(\alpha^{-1})
  =\operatorname N(\alpha)^{-1}=-\frac1{11}
\]
이고
\[
  \operatorname N(\alpha)\operatorname N(\alpha^{-1})=1
\]
이므로 검산된다.
\end{solution}

\subsection*{연습문제 \ref{ex:norm-trace-pure-cubic}}

\begin{solution}
$\theta$의 최소다항식은
\[
  f(T)=T^3-2
\]
이다. 따라서 계수 공식에서
\[
  \Tr(\theta)=0,
  \qquad
  \operatorname N(\theta)=(-1)^3(-2)=2.
\]
노름의 곱셈성으로
\[
  \operatorname N(\theta^2)=\operatorname N(\theta)^2=4.
\]
대각합은 켤레들의 합으로 계산할 수 있다. $\theta$의 켤레가
\[
  \theta,\quad\zeta_3\theta,\quad\zeta_3^2\theta
\]
이므로
\[
  \Tr(\theta^2)
  =\theta^2(1+\zeta_3^2+\zeta_3^4)
  =\theta^2(1+\zeta_3+\zeta_3^2)=0.
\]
또는 뉴턴 항등식에서 바로 $s_2=0$을 얻는다.

대각합의 선형성에서
\[
  \Tr(1+\theta)=\Tr(1)+\Tr(\theta)=3+0=3.
\]
노름은 정리 \ref{thm:norm-resultant}를 사용한다.
\[
\begin{aligned}
  \operatorname N(1+\theta)
  &=\operatorname N(\theta-(-1))\\
  &=(-1)^3f(-1)\\
  &=-((-1)^3-2)=3.
\end{aligned}
\]
마지막으로
\[
  \operatorname N(2-\theta)
  =\prod_i(2-\theta_i)=f(2)=8-2=6.
\]
여기서 $2-\theta_i$의 곱은 최소다항식의 값 $f(2)$과 같은 방향의 차이를 사용하므로 추가 부호가 없다.
\end{solution}

\subsection*{연습문제 \ref{ex:characteristic-polynomial-repeated}}

\begin{solution}
$\alpha=\sqrt2$의 $\Q$ 위 최소다항식은
\[
  f_\alpha(T)=T^2-2
\]
이고
\[
  [L:\Q(\alpha)]
  =[\Q(\sqrt2,\sqrt3):\Q(\sqrt2)]=2.
\]
정리 \ref{thm:characteristic-minimal-power}에 의해
\[
  P_{\alpha,L/\Q}(T)
  =(T^2-2)^2=T^4-4T^2+4.
\]
따라서 $T^3$의 계수가 $0$이고 상수항이 $4$이므로
\[
  \Tr_{L/\Q}(\sqrt2)=0,
  \qquad
  \operatorname N_{L/\Q}(\sqrt2)=4.
\]

한편 $E=\Q(\sqrt2)$에서는
\[
  \Tr_{E/\Q}(\sqrt2)=0,
  \qquad
  \operatorname N_{E/\Q}(\sqrt2)=-2.
\]
$\sqrt2\in E$이고 $[L:E]=2$이므로
\[
  \Tr_{L/\Q}(\sqrt2)
  =2\Tr_{E/\Q}(\sqrt2)=0,
\]
\[
  \operatorname N_{L/\Q}(\sqrt2)
  =\operatorname N_{E/\Q}(\sqrt2)^2=4.
\]
전체 확장에서 각 $E/\Q$-켤레가 두 번씩 반복되는 것이 차이의 원인이다.
\end{solution}

\subsection*{연습문제 \ref{ex:finite-field-f8-trace-norm}}

\begin{solution}
$\F_8/\F_2$의 갈루아군은 $x\mapsto x^2$로 생성되므로
\[
  \Tr(x)=x+x^2+x^4.
\]
먼저 표수 $2$에서
\[
  \Tr(1)=1+1+1=1.
\]
관계 $\theta^3+\theta+1=0$에서
\[
  \theta^3=\theta+1,
  \qquad
  \theta^4=\theta^2+\theta
\]
이므로
\[
  \Tr(\theta)
  =\theta+\theta^2+(\theta^2+\theta)=0.
\]
또한 $\theta^8=\theta$이므로
\[
\begin{aligned}
  \Tr(\theta^2)
  &=\theta^2+\theta^4+\theta^8\\
  &=\theta^2+(\theta^2+\theta)+\theta=0.
\end{aligned}
\]
대각합의 $\F_2$-선형성에서
\[
  \Tr(a+b\theta+c\theta^2)=a.
\]
따라서
\[
  \ker\Tr
  =\{0,\theta,\theta^2,\theta+\theta^2\}.
\]

노름은
\[
  \operatorname N(x)=x^{1+2+4}=x^7.
\]
$\F_8^\times$의 위수가 $7$이므로 모든 $x\ne0$에 대해 $x^7=1$이다. 따라서
\[
  \operatorname N(x)=1
  \qquad(x\in\F_8^\times).
\]
\end{solution}

\subsection*{연습문제 \ref{ex:norm-trace-tower-sqrt2-sqrt3}}

\begin{solution}
$L/M$는 $\sqrt3$의 부호를 바꾸는 자동동형
\[
  \tau:\sqrt3\mapsto-\sqrt3,
  \qquad \sqrt2\mapsto\sqrt2
\]
를 갖는 이차 갈루아 확장이다. 따라서
\[
\begin{aligned}
  \Tr_{L/M}(\alpha)
  &=\alpha+\tau(\alpha)\\
  &=(\sqrt2+\sqrt3)+(\sqrt2-\sqrt3)\\
  &=2\sqrt2,
\end{aligned}
\]
\[
\begin{aligned}
  \operatorname N_{L/M}(\alpha)
  &=\alpha\tau(\alpha)\\
  &=(\sqrt2+\sqrt3)(\sqrt2-\sqrt3)\\
  &=2-3=-1.
\end{aligned}
\]

추이성을 적용하면
\[
  \Tr_{L/\Q}(\alpha)
  =\Tr_{M/\Q}(2\sqrt2)=0
\]
이고
\[
  \operatorname N_{L/\Q}(\alpha)
  =\operatorname N_{M/\Q}(-1)=(-1)^2=1.
\]

직접 검산하면 네 켤레는
\[
  \sqrt2+\sqrt3,
  \quad\sqrt2-\sqrt3,
  \quad-\sqrt2+\sqrt3,
  \quad-\sqrt2-\sqrt3.
\]
합은 $0$이고, 첫 두 항의 곱과 마지막 두 항의 곱이 각각 $-1$이므로 전체 곱은 $1$이다.
\end{solution}

\subsection*{연습문제 \ref{ex:cyclotomic-zeta5-norm-trace}}

\begin{solution}
$\zeta=\zeta_5$의 최소다항식은
\[
  \Phi_5(T)=T^4+T^3+T^2+T+1
\]
이다. 계수 공식에서
\[
  \Tr_{L/\Q}(\zeta)=-1,
  \qquad
  \operatorname N_{L/\Q}(\zeta)=1.
\]
또는 네 켤레 $\zeta,\zeta^2,\zeta^3,\zeta^4$를 직접 합하고 곱해도 된다.

대각합의 선형성에서
\[
\begin{aligned}
  \Tr(1-\zeta)
  &=\Tr(1)-\Tr(\zeta)\\
  &=4-(-1)=5.
\end{aligned}
\]
노름은
\[
\begin{aligned}
  \operatorname N(1-\zeta)
  &=\prod_{a=1}^4(1-\zeta^a)\\
  &=\Phi_5(1)=5.
\end{aligned}
\]
대각합과 노름이 우연히 같은 값 $5$가 되지만, 두 계산의 원리는 각각 합과 곱으로 다르다.
\end{solution}

\subsection*{연습문제 \ref{ex:prove-artin-independence}}

\begin{solution}
비자명한 관계가 존재한다고 가정하고 항의 수가 최소인 관계
\[
  a_1\chi_1+\cdots+a_r\chi_r=0
\]
를 택한다. 영인 계수의 항은 제거할 수 있으므로 모든 $a_i\ne0$이다. $r=1$이면 항등원에서 평가하여 $a_1=0$이 되므로 $r\ge2$이다.

$\chi_1\ne\chi_r$이므로 어떤 $g\in G$에 대해
\[
  \chi_1(g)-\chi_r(g)\ne0
\]
이다. 임의의 $h\in G$에 대해 관계를 $gh$에 적용하면
\[
  \sum_{i=1}^r a_i\chi_i(g)\chi_i(h)=0.
\]
관계를 $h$에 적용한 뒤 $\chi_r(g)$를 곱한 식을 빼면
\[
  \sum_{i=1}^{r-1}
  a_i(\chi_i(g)-\chi_r(g))\chi_i(h)=0
\]
이다. $i=r$인 항은 정확히 사라진다. 한편 $i=1$인 계수
\[
  a_1(\chi_1(g)-\chi_r(g))
\]
는 두 인수가 모두 영이 아니므로 영이 아니다. 따라서 이 식은 $r-1$개 이하의 캐릭터 사이의 비자명한 관계이다. 이는 $r$의 최소성에 모순이다. 그러므로 처음의 비자명한 관계는 존재하지 않는다.
\end{solution}

\subsection*{연습문제 \ref{ex:prove-embedding-formulas}}

\begin{solution}
$E=K(\alpha)$, $d=[E:K]$, $r=[L:E]$라 하자. $L/K$가 분리이므로 $E/K$와 $L/E$도 분리이다. $\alpha$의 최소다항식의 서로 다른 근들을
\[
  \alpha_1,\dots,\alpha_d
\]
라 쓰면 $E$의 $K$-매장들은 $\alpha$를 각각 $\alpha_j$로 보낸다.

분리확장에서 매장의 연장 개수는 상대차수와 같다. 따라서 $E$의 각 $K$-매장은 $L$의 $K$-매장으로 정확히
\[
  [L:E]=r
\]
가지 방식으로 연장된다. 그러므로 $L$의 모든 $K$-매장 $\sigma_i$를 움직일 때 값 $\sigma_i(\alpha)$의 목록에는 각 $\alpha_j$가 정확히 $r$번 나타난다. 따라서
\[
\begin{aligned}
  \prod_{i=1}^{dr}(T-\sigma_i(\alpha))
  &=\prod_{j=1}^d(T-\alpha_j)^r\\
  &=f_\alpha(T)^r.
\end{aligned}
\]
정리 \ref{thm:characteristic-minimal-power}에서
\[
  P_{\alpha,L/K}=f_\alpha^r
\]
이므로
\[
  P_{\alpha,L/K}(T)
  =\prod_i(T-\sigma_i(\alpha)).
\]
둘째 계수와 상수항을 비교하면
\[
  \Tr_{L/K}(\alpha)=\sum_i\sigma_i(\alpha),
  \qquad
  \operatorname N_{L/K}(\alpha)=\prod_i\sigma_i(\alpha)
\]
를 얻는다.
\end{solution}
